\section{Learning corrections outside the frozen model}
\label{sec:adapt}
Consider again the weakened motor. We compare the model's predicted outcome
with the outcome observed in the environment, then use the difference to
update two small objects. A score vector identifies a suitable command
correction; a residual matrix improves the predicted outcome. The following description uses
the decision index $n$ and target $d_n$ defined in Equation~\eqref{eq:action}.

\subsection{Predict first, update after observation}
At deployment, the selected predictor parameters $\theta_0$, semantic heads,
mechanism router, and underlying control policy are frozen. The online interface
receives the current observation, executed command, task identifier, and eventual
outcome. Reward, success, and the true actuator change are evaluation variables,
not online update inputs.

Every prediction is recorded before its target is observed. After the horizon
has elapsed, the system computes an error and updates external state for later
decisions. The first four transitions in an episode supply evidence to the
frozen mechanism router. The route locks after the fourth outcome arrives. The system can then update
the persistent state for the first time. Consequently, decision five is
the earliest command that can change. Appendix~\ref{app:algorithm} gives the
complete schedule, including the different lifetimes of episode prefixes and
persistent matrices.

\subsection{Scores select an actuator correction}
The model compares a bank of $J=32$ transformations: identity, shear, and
attenuation candidates. Identity leaves the command unchanged. Shear couples
the first two action coordinates, so that an $x$ command also affects $y$.
Attenuation reduces the realized movement in the $x/y$ coordinates. The bank's
values are registered offline; the online system infers a candidate from
observed transitions.

For candidate $j$, let $T_j$ be its command transformation. Once outcome $d_n$
is available, the candidate receives loss
\begin{equation}
 \ell_{n,j}=\frac{1}{D_q}
 \|g_{\theta_0}(q_n,T_j(a_n),c)-d_n\|_2^2,
 \qquad S_j\leftarrow S_j+\ell_{n,j}.
 \label{eq:scores}
\end{equation}
Here $S_j$ is its cumulative score. A lower loss means that the candidate transformation better explains the
observed transition through the frozen predictor.
A frozen router first selects the mechanism family $\rho$ from the four
initial transitions; within that family, the candidate with the smallest
$S_j$ supplies the estimated transformation. These comparisons always query
the original frozen predictor, without the residual matrix described below.

Let $a_n^*$ be the command proposed by the underlying policy, and let
$\tilde a_n$ be the selected inverse transformation applied to that command.
The controller executes
\begin{equation}
 a_n=a_n^*+\operatorname{clip}(\tilde a_n-a_n^*,-\delta,\delta),
 \qquad \delta=0.05.
 \label{eq:command}
\end{equation}
The clip acts separately on each coordinate and caps the size of the correction.
If this capped command would lie outside the allowed action range, the system
uses $a_n^*$ instead. The inverse maps, gain floor, and source/noise behavior
are given in Appendix~\ref{app:algorithm}.

For a one-coordinate illustration, suppose the route has locked and prior
observations have selected gain 0.5. A desired command of 0.10 gives inverse
proposal 0.20. The 0.05 cap makes the executed command 0.15. If the actual gain
is 0.5, its realized effect is 0.075, compared with 0.05 without correction.
Candidate scoring uses the forward gain transform to explain observed outcomes;
the inverse compensates a later command. The residual matrix below changes the
predicted outcome, not this command. These illustrative values explain the
operation rather than report an additional experiment.

\subsection{Matrices refine the prediction}
After selecting a candidate, the frozen core computes hidden feature
$h_n\in\R^d$ using its transformed command. We retain the fixed offline mean
$\bar h\in\R^d$ and projection $B\in\R^{r\times d}$, where $r=30$.
The projection is obtained by principal component analysis of offline action
features. Each correctable mechanism family has its own matrix
$C_\rho\in\R^{d\times r}$, initialized to zero. Prediction becomes
\begin{equation}
 z_n=B(h_n-\bar h),\qquad
 \tilde h_n=h_n+C_\rho z_n,\qquad
 \hat d_n=W_{\mathrm{out}}\tilde h_n+b_{\mathrm{out}}.
 \label{eq:correction}
\end{equation}
The vector $z_n\in\R^r$ compresses the centered feature. Multiplication by
$C_\rho$ expands it into a learned correction to the hidden representation.
The original readout has weight $W_{\mathrm{out}}\in\R^{D_q\times d}$ and
bias $b_{\mathrm{out}}\in\R^{D_q}$; both remain frozen.
The prediction changes through the residual matrix; the feature extractor
and readout keep their trained weights.

After observing an outcome and updating the candidate scores, the active
matrix takes a normalized gradient step on prediction MSE. The step length is
$\eta=0.02$, and the matrix is projected to a Frobenius-norm ball of radius
$R_C=0.20$. Appendix~\ref{app:matrix-update} spells out this update and explains
how it differs from the earlier pre-update error used for evaluation. Candidate scoring and the action inverse do not read the matrices. We can
therefore remove the matrices to measure their prediction effect while
leaving action selection unchanged.

\subsection{What remains fixed, and what occupies memory?}
The trained semantic computation has no read connection to the external
matrices or scores. Given the same sensory input and inference settings,
it therefore returns the same logits (pre-softmax class scores) before and
after online adaptation.
We verify this with parameter and fixed-input-output hashes. This property
preserves the learned mapping; semantic outputs can still differ when the
robot observes different inputs.

The online allocation contains two residual matrices, the score vector and
count, compatibility padding, and a four-transition routing prefix. Their peak
total is 28,848 bytes of mutable numeric state. This is the memory that can
change during adaptation, separate from the frozen model and transient
inference activations. Table~\ref{tab:ledger} gives the exact arithmetic.
The bounded prefix is cleared after routing, and the matrices and scores retain
experience across the stream.
